% 来源：OEBPS/Text/01_41.xhtml（OPF spine 顺序；仅第一册）
\cleardoublepage
\BookEntry[狄更斯《圣诞欢歌》]{狄更斯《圣诞欢歌》}

\BookCover{images/书单25圣诞欢歌.jpg}{[英] 查理·狄更斯著| 王漪译| 二十一世纪出版社2010年版}

\begin{BookSummaryBox}
\textbf{狄更斯最深入人心的一部小说}
\par
\end{BookSummaryBox}

\subsubsection*{这本书谈什么？}

主人公斯克鲁济是一位商人，一生只知道攒钱，毫无人性。他克扣员工的工资，拒绝外甥的问候和祝福，不施舍分文给穷人。在平安夜，三个精灵带领他去领略了过去、现在和未来的一些生活情境，让斯克鲁济从自制的陷阱和枷锁中解脱出来，领悟到圣诞节的意义和生活的意义。他在圣诞颂歌奏响时开始了自己新的人生。
\par

这本书让我们意识到，圣诞节不是一道吃鹅还是吃火鸡的客观题，而是一道反省人与人之间——既然都是神的创造——应当如何互相关爱的主观题。《圣诞颂歌》中的人物和情节在西方家喻户晓，一百多年来积聚的人气使其成为狄更斯小说中最深入人心的一部。社会学家相信，这部作品深刻地改变了现代圣诞节在西方世界的地位和价值，狄更斯本人因此被称为“发明圣诞节”的人。
\par

\subsubsection*{芒格为什么谈起这本书？}

芒格在关于“人类误判心理学”的演讲中，说到要“避免不一致倾向”：
\par

\begin{quotation}
为了节省运算空间，人类的大脑会不愿意作出改变。这是一种避免不一致性的形式。在所有的人类习惯中，无论是好习惯还是坏习惯，我们都能看到这种情况。没几个人能够列出许多他们已经改掉的坏习惯，而有些人哪怕连一个都列举不出来。与此相反，几乎每个人都有大量持续很久的坏习惯，尽管他们自己也知道这些习惯不好。考虑到这种情况，在许多时候我们说“三岁看老”不是没有道理的。在狄更斯的《圣诞欢歌》中，可怜的雅各布·马里的鬼魂说：“我戴着我在生活中锻造的锁链。”他说的锁链就是那些起初轻微得难以察觉，在察觉之后又牢固得无法打破的习惯。
\par
\end{quotation}

\subsubsection*{《圣诞欢歌》作者狄更斯如是说}

“‘斯克鲁济，你太害怕这个世界了，’姑娘温和地说，‘你把一切其他希望都并入了这个追求，就是为了避免这个世界上肮脏的责备。我已经看见你原来高尚的志向都一个紧接一个地消失了，只剩下那个主要的欲望，那就是唯利是图。’”
\par

“‘这些悲惨的孩子属于人类。’幽灵说。‘男孩叫罪恶，女孩叫需要。如果不为他们做点什么，他们将毁掉人类。’”
\par

“‘我戴着我在生活中锻造的锁链。’雅各布·马利的鬼魂说。”
\par
